% GENERATED FILE. Edit canonical lessons, then run npm run generate:books.
% canonical-source-hash: fnv1a64:9f789eac619412d2
% canonical-lessons: MW-C256-darno, MW-C256-dukh, MW-C256-sukh, MW-C256-anand, MW-C256-umang

\chapter{Fear and Joy}
\label{ch:mw-fear-and-joy}
\hlchaptermodality{256}

\begin{chapteropening}
\textbf{By the end of this chapter:} \emph{I can say what frightens me and what makes me glad.}

Complete the last lesson of chapter 256: I can say what frightens me and what makes me glad.
\end{chapteropening}

\section[\texorpdfstring{\mw{डरणो} (ḍarṇo)}{डरणो (ḍarṇo)}]{\mw{डरणो} (ḍarṇo) — to be afraid}

\label{lesson:MW-C256-darno}



\begin{quote}
Before the new one: say the Marwadi for a tear (from the eye), then the Marwadi for to tremble.
\end{quote}

\subsection*{You'll want to know: \mw{डरणो}}
\textbf{\mw{डरणो}} — \emph{ḍarṇo} — \textquotedblleft{}to be afraid\textquotedblright{}.

\begin{quote}
\mw{छोरी} \mw{सांप} \mw{सूं} \mw{डरै} \mw{है।}
\end{quote}

\begin{quote}
\emph{chhorī sā̃p sū̃ ḍarai hai.}
\end{quote}

\begin{quote}
The girl is afraid of snakes.
\end{quote}

Past (he …): \mw{वो} \mw{छोरो} \mw{डर्यो}, \emph{vo chhoro ḍaryo}.

\begin{cousinweb}
\textbf{\mw{डरणो}} (\emph{ḍarṇo}) is inherited from Sanskrit \emph{darati}, \textquotedblleft{}is frightened\textquotedblright{}; CDIAL lists Marwari \emph{ḍarṇo} under it (CDIAL 6190).
\end{cousinweb}

\begin{grammarlens}[title={What you've built}]
One more word for this chapter.
\end{grammarlens}

\subsection*{Your turn}
\begin{itemize}
\raggedright
  \item \emph{Say it:} \emph{ḍarṇo}
  \item \emph{Say it:} \emph{ḍarṇo}, once more
  \item \emph{Say it:} \emph{chhorī sā̃p sū̃ ḍarai hai.} — The girl is afraid of snakes.
  \item \emph{Recall:} say the Marwadi for a tear (from the eye), then the Marwadi for to tremble, then say \emph{ḍarṇo} again
\end{itemize}

\begin{tcolorbox}[breakable,colback=teal!4,colframe=teal!35!black,title={Before you move on}]
What does \mw{डरणो} mean? (\textquotedblleft{}To be afraid\textquotedblright{}.) Say it once more.
\end{tcolorbox}
\section[\texorpdfstring{\mw{दुख} (dukh)}{दुख (dukh)}]{\mw{दुख} (dukh) — sorrow; grief}

\label{lesson:MW-C256-dukh}



\begin{quote}
Before the new one: say the Marwadi for to tremble, then the Marwadi for to be afraid.
\end{quote}

\subsection*{You'll want to know: \mw{दुख}}
\textbf{\mw{दुख}} — \emph{dukh} — \textquotedblleft{}sorrow; grief\textquotedblright{}.

\begin{quote}
\mw{गाय} \mw{मरी}, \mw{घर} \mw{मांय} \mw{दुख} \mw{है।}
\end{quote}

\begin{quote}
\emph{gāy marī, ghar mā̃y dukh hai.}
\end{quote}

\begin{quote}
The cow died; there is sorrow in the house.
\end{quote}

A masculine noun.

\begin{cousinweb}
\textbf{\mw{दुख}} (\emph{dukh}) is inherited from Sanskrit \emph{duḥkha}, \textquotedblleft{}difficulty, pain\textquotedblright{}; CDIAL lists Marwari \emph{dukh}, \textquotedblleft{}sorrow\textquotedblright{}, under it (CDIAL 6375).
\end{cousinweb}

\begin{grammarlens}[title={What you've built}]
One more word for this chapter.
\end{grammarlens}

\subsection*{Your turn}
\begin{itemize}
\raggedright
  \item \emph{Say it:} \emph{dukh}
  \item \emph{Say it:} \emph{dukh}, once more
  \item \emph{Say it:} \emph{gāy marī, ghar mā̃y dukh hai.} — The cow died; there is sorrow in the house.
  \item \emph{Recall:} say the Marwadi for to tremble, then the Marwadi for to be afraid, then say \emph{dukh} again
\end{itemize}

\begin{tcolorbox}[breakable,colback=teal!4,colframe=teal!35!black,title={Before you move on}]
What does \mw{दुख} mean? (\textquotedblleft{}Sorrow; grief\textquotedblright{}.) Say it once more.
\end{tcolorbox}
\section[\texorpdfstring{\mw{सुख} (sukh)}{सुख (sukh)}]{\mw{सुख} (sukh) — ease and contentment; happiness}

\label{lesson:MW-C256-sukh}



\begin{quote}
Before the new one: say the Marwadi for to be afraid, then the Marwadi for sorrow.
\end{quote}

\subsection*{You'll want to know: \mw{सुख}}
\textbf{\mw{सुख}} — \emph{sukh} — \textquotedblleft{}ease and contentment; happiness\textquotedblright{}.

\begin{quote}
\mw{गांव} \mw{मांय} \mw{सुख} \mw{है।}
\end{quote}

\begin{quote}
\emph{gā̃v mā̃y sukh hai.}
\end{quote}

\begin{quote}
Life in the village is content.
\end{quote}

A masculine noun.

\begin{cousinweb}
\textbf{\mw{सुख}} (\emph{sukh}) is a tatsama, the Sanskrit \emph{sukha}, \textquotedblleft{}ease, comfort, happiness\textquotedblright{}, taken over with its final vowel dropped in speech (Monier-Williams).
\end{cousinweb}

\begin{grammarlens}[title={What you've built}]
One more word for this chapter.
\end{grammarlens}

\subsection*{Your turn}
\begin{itemize}
\raggedright
  \item \emph{Say it:} \emph{sukh}
  \item \emph{Say it:} \emph{sukh}, once more
  \item \emph{Say it:} \emph{gā̃v mā̃y sukh hai.} — Life in the village is content.
  \item \emph{Recall:} say the Marwadi for to be afraid, then the Marwadi for sorrow, then say \emph{sukh} again
\end{itemize}

\begin{tcolorbox}[breakable,colback=teal!4,colframe=teal!35!black,title={Before you move on}]
What does \mw{सुख} mean? (\textquotedblleft{}Ease and contentment; happiness\textquotedblright{}.) Say it once more.
\end{tcolorbox}
\section[\texorpdfstring{\mw{आणंद} (āṇand)}{आणंद (āṇand)}]{\mw{आणंद} (āṇand) — joy; delight}

\label{lesson:MW-C256-anand}



\begin{quote}
Before the new one: say the Marwadi for sorrow, then the Marwadi for contentment.
\end{quote}

\subsection*{You'll want to know: \mw{आणंद}}
\textbf{\mw{आणंद}} — \emph{āṇand} — \textquotedblleft{}joy; delight\textquotedblright{}.

\begin{quote}
\mw{ब्याव} \mw{मांय} \mw{घणो} \mw{आणंद} \mw{है।}
\end{quote}

\begin{quote}
\emph{byāv mā̃y ghaṇo āṇand hai.}
\end{quote}

\begin{quote}
There is great joy at the wedding.
\end{quote}

A masculine noun.

\begin{cousinweb}
\textbf{\mw{आणंद}} (\emph{āṇand}) is inherited from Sanskrit \emph{ānanda}, \textquotedblleft{}happiness\textquotedblright{}. CDIAL gives no Marwari form in this entry; it lists Gujarati \emph{āṇãd}, \textquotedblleft{}joy\textquotedblright{}, with the same \mw{ण} (CDIAL 1172).
\end{cousinweb}

\begin{grammarlens}[title={What you've built}]
One more word for this chapter.
\end{grammarlens}

\subsection*{Your turn}
\begin{itemize}
\raggedright
  \item \emph{Say it:} \emph{āṇand}
  \item \emph{Say it:} \emph{āṇand}, once more
  \item \emph{Say it:} \emph{byāv mā̃y ghaṇo āṇand hai.} — There is great joy at the wedding.
  \item \emph{Recall:} say the Marwadi for sorrow, then the Marwadi for contentment, then say \emph{āṇand} again
\end{itemize}

\begin{tcolorbox}[breakable,colback=teal!4,colframe=teal!35!black,title={Before you move on}]
What does \mw{आणंद} mean? (\textquotedblleft{}Joy; delight\textquotedblright{}.) Say it once more.
\end{tcolorbox}
\section[\texorpdfstring{\mw{उमंग} (umaṅg)}{उमंग (umaṅg)}]{\mw{उमंग} (umaṅg) — high spirits; eagerness}

\label{lesson:MW-C256-umang}



\begin{quote}
Before the new one: say the Marwadi for contentment, then the Marwadi for joy.
\end{quote}

\subsection*{You'll want to know: \mw{उमंग}}
\textbf{\mw{उमंग}} — \emph{umaṅg} — \textquotedblleft{}high spirits; eagerness\textquotedblright{}.

\begin{quote}
\mw{आज} \mw{मन} \mw{मांय} \mw{घणी} \mw{उमंग} \mw{है।}
\end{quote}

\begin{quote}
\emph{āj man mā̃y ghaṇī umaṅg hai.}
\end{quote}

\begin{quote}
Today my spirits are high.
\end{quote}

A feminine noun.

\begin{cousinweb}
\textbf{\mw{उमंग}} (\emph{umaṅg}) is inherited from Sanskrit \emph{unmajjati}, \textquotedblleft{}emerges\textquotedblright{}: a feeling that wells up. CDIAL gives no Marwari form in this entry; it lists Hindi \emph{umãg}, \textquotedblleft{}elation\textquotedblright{} (CDIAL 2110).
\end{cousinweb}

\begin{grammarlens}[title={What you've built}]
One more word for this chapter.
\end{grammarlens}

\subsection*{Your turn}
\begin{itemize}
\raggedright
  \item \emph{Say it:} \emph{umaṅg}
  \item \emph{Say it:} \emph{umaṅg}, once more
  \item \emph{Say it:} \emph{āj man mā̃y ghaṇī umaṅg hai.} — Today my spirits are high.
  \item \emph{Recall:} say the Marwadi for contentment, then the Marwadi for joy, then say \emph{umaṅg} again
\end{itemize}

\begin{tcolorbox}[breakable,colback=teal!4,colframe=teal!35!black,title={Before you move on}]
What does \mw{उमंग} mean? (\textquotedblleft{}High spirits; eagerness\textquotedblright{}.) Say it once more.
\end{tcolorbox}
